%========================================
% MatinBook - Stage 15: Cover Test
% Test: Front Cover + Back Cover + Full Book Integration
% Purpose: Validate cover design and complete book assembly
%========================================

% Add parent directory to search path
\makeatletter
\def\input@path{{../}}
\makeatother

\documentclass{matinbook}

\title{هنر برنامه‌نویسی}
\author{تیم توسعه MatinBook}
\date{تابستان ۱۴۰۵}

\begin{document}
    
    %========================================
    % FRONT COVER
    %========================================
    \makecover
    {هنر برنامه‌نویسی}
    {از ایده تا اجرا}
    {تیم توسعه MatinBook}
    {تابستان ۱۴۰۵}
    
    %========================================
    % BACK COVER
    %========================================
    \makebackcover{%
        \textbf{هنر برنامه‌نویسی} کتابی جامع برای
        یادگیری اصول برنامه‌نویسی و الگوریتم است.
        
        \textbf{در این کتاب می‌آموزید:}
        \begin{itemize}
            \item مفاهیم پایه برنامه‌نویسی با پایتون
            \item الگوریتم‌های جستجو و مرتب‌سازی
            \item ساختمان داده‌های پرکاربرد
            \item تحلیل پیچیدگی و بهینه‌سازی
            \item برنامه‌نویسی شیءگرا
        \end{itemize}
        
        \textbf{ویژگی‌های کتاب:}
        \begin{itemize}
            \item بیش از ۵۰ مثال کاربردی
            \item تمرین‌های متنوع با پاسخ تشریحی
            \item نمودارها و شبه‌کدهای استاندارد
            \item نمایه جامع واژگان
        \end{itemize}
        
        \vspace{0.5cm}
        \textbf{ساخته شده با \lr{MatinBook} نسخه ۱}
    }
    
    %========================================
    % TITLE PAGE
    %========================================
    \maketitle
    
    %========================================
    % COPYRIGHT
    %========================================
    \thispagestyle{empty}
    \vspace*{5cm}
    \begin{center}
        {\large\textbf{حقوق نشر}}\\[1cm]
        تمامی حقوق این کتاب محفوظ است.\\[0.5cm]
        نسخه ۱.۰ --- تابستان ۱۴۰۵\\[1cm]
        \lr{Copyright \copyright\ 2026 MatinBook Project}\\
        \lr{Licensed under MIT License}\\[0.5cm]
        \textit{این کتاب با استفاده از کلاس \lr{MatinBook} نسخه ۱}
        \textit{حروف‌چینی شده است.}
    \end{center}
    \clearpage
    
    %========================================
    % DEDICATION
    %========================================
    \thispagestyle{empty}
    \vfill
    \begin{flushright}
        {\Large\itshape
            به تمام کسانی که\\
            با شوق می‌آموزند\\
            و با عشق می‌آموزانند
        }
    \end{flushright}
    \clearpage
    \frontmatter
    
    %========================================
    % PREFACE
    %========================================
    \chapter*{پیشگفتار}
    \addcontentsline{toc}{chapter}{پیشگفتار}
    
    این کتاب برای کسانی نوشته شده که می‌خواهند
    برنامه‌نویسی را به صورت اصولی بیاموزند.
    مطالب از سطح پایه شروع شده و به تدریج
    به مباحث پیشرفته می‌رسد.
    
    امیدواریم این کتاب بتواند سهمی هرچند کوچک
    در ارتقای دانش برنامه‌نویسی فارسی‌زبانان
    داشته باشد.
    
    \vspace{1cm}
    \begin{flushleft}
        \textbf{تیم توسعه MatinBook}\\
        تابستان ۱۴۰۵
    \end{flushleft}
    
    \clearpage
    
    %========================================
    % TABLE OF CONTENTS
    %========================================
    \tableofcontents
    \clearpage
    \mainmatter
    
    %========================================
    % Chapter 1: Introduction
    %========================================
    \chapter{مقدمه}
    
    \section{چرا برنامه‌نویسی؟}
    
    برنامه‌نویسی\index{برنامه‌نویسی} مهارتی است
    که در دنیای امروز به اندازه خواندن و نوشتن
    اهمیت دارد. از تلفن‌های همراه گرفته
    تا خودروهای خودران، همه با کد کار می‌کنند.
    
    \begin{definition}[برنامه‌نویسی]
        برنامه‌نویسی فرایند طراحی، نوشتن،
        آزمایش و نگهداری کدهای کامپیوتری است
        که به ماشین می‌گویند چه کاری انجام دهد.
        \label{def:programming}
    \end{definition}
    
    \section{انتخاب زبان برنامه‌نویسی}
    
    برای شروع، پایتون\index{پایتون}
    بهترین انتخاب است. دلایل:
    
    \begin{itemize}
        \item نحو ساده و خوانا
        \item جامعه کاربری بزرگ
        \item کتابخانه‌های فراوان
        \item کاربرد در حوزه‌های مختلف
    \end{itemize}
    
    \begin{example}[اولین برنامه]
        ساده‌ترین برنامه در پایتون:
        \begin{latin}
            \begin{minted}{python}
                print("Hello, World!")
            \end{minted}
        \end{latin}
        \label{ex:hello-world}
    \end{example}
    
    %========================================
    % Chapter 2: Basics
    %========================================
    \chapter{مفاهیم پایه}
    
    \section{متغیرها}
    
    متغیر\index{متغیر} محلی برای ذخیره داده است.
    در پایتون، متغیرها نیاز به اعلان نوع ندارند:
    
    \begin{latin}
        \begin{minted}{python}
            # Variables in Python
            name = "Ali"           # String
            age = 25               # Integer
            pi = 3.14              # Float
            is_active = True       # Boolean
            
            print(f"{name} is {age} years old")
        \end{minted}
    \end{latin}
    \codecaption{تعریف متغیر در پایتون}
    \label{code:variables}
    
    \section{عملگرها}
    
    عملگرها\index{عملگر} عملیات روی داده‌ها
    را انجام می‌دهند:
    
    \begin{table}[h]
        \centering
        \caption{عملگرهای پایه در پایتون}
        \label{tab:operators}
        \begin{tabular}{|c|c|c|}
            \hline
            \textbf{نوع} & \textbf{عملگر} & \textbf{مثال} \\
            \hline
            محاسباتی & \lr{+, -, *, /, //, \%, **} & \lr{5 + 3 = 8} \\
            مقایسه‌ای & \lr{==, !=, <, >, <=, >=} & \lr{5 > 3 = True} \\
            منطقی & \lr{and, or, not} & \lr{True and False = False} \\
            \hline
        \end{tabular}
    \end{table}
    
    %========================================
    % Chapter 3: Control Flow
    %========================================
    \chapter{کنترل جریان}
    
    \section{دستورات شرطی}
    
    \begin{latin}
        \begin{minted}{python}
            score = 85
            
            if score >= 90:
            grade = "A"
            elif score >= 80:
            grade = "B"
            elif score >= 70:
            grade = "C"
            else:
            grade = "F"
            
            print(f"Score: {score}, Grade: {grade}")
        \end{minted}
    \end{latin}
    \codecaption{دستور شرطی در پایتون}
    \label{code:if-else}
    
    \section{حلقه‌ها}
    
    \begin{latin}
        \begin{minted}{python}
            # For loop
            for i in range(5):
            print(f"Iteration {i}")
            
            # While loop
            count = 0
            while count < 5:
            print(f"Count: {count}")
            count += 1
        \end{minted}
    \end{latin}
    \codecaption{حلقه‌های \lr{for} و \lr{while}}
    \label{code:loops}
    
    %========================================
    % Chapter 4: Functions
    %========================================
    \chapter{توابع}
    
    \section{تعریف و استفاده}
    
    تابع\index{تابع} بلوکی از کد است که
    یک وظیفه مشخص را انجام می‌دهد:
    
    \begin{latin}
        \begin{minted}{python}
            def greet(name: str) -> str:
            """Return a greeting message."""
            return f"Hello, {name}!"
            
            def add(a: int, b: int) -> int:
            """Return the sum of two numbers."""
            return a + b
            
            # Using functions
            print(greet("Ali"))
            print(f"3 + 4 = {add(3, 4)}")
        \end{minted}
    \end{latin}
    \codecaption{تعریف و استفاده از توابع}
    \label{code:functions}
    
    \begin{exercise}
        تابعی بنویسید که یک عدد دریافت کند
        و فاکتوریل آن را محاسبه کند.
        \label{exc:factorial}
    \end{exercise}
    
    \begin{solution}
        \begin{latin}
            \begin{minted}{python}
                def factorial(n: int) -> int:
                """Calculate factorial recursively."""
                if n <= 1:
                return 1
                return n * factorial(n - 1)
                
                print(factorial(5))  # 120
            \end{minted}
        \end{latin}
    \end{solution}
    
    %========================================
    % Chapter 5: Data Structures
    %========================================
    \chapter{ساختمان داده‌ها}
    
    \section{لیست}
    
    لیست\index{لیست} پرکاربردترین ساختمان داده
    در پایتون است:
    
    \begin{latin}
        \begin{minted}{python}
            # Create and manipulate lists
            numbers = [1, 2, 3, 4, 5]
            numbers.append(6)        # Add to end
            numbers.insert(0, 0)     # Insert at index
            numbers.remove(3)        # Remove by value
            
            print(f"List: {numbers}")
            print(f"Length: {len(numbers)}")
            print(f"Sum: {sum(numbers)}")
        \end{minted}
    \end{latin}
    \codecaption{کار با لیست در پایتون}
    \label{code:lists}
    
    \section{دیکشنری}
    
    دیکشنری\index{دیکشنری} داده‌ها را به صورت
    کلید-مقدار\index{کلید-مقدار} ذخیره می‌کند:
    
    \begin{latin}
        \begin{minted}{python}
            # Create a dictionary
            student = {
                "name": "Ali",
                "age": 25,
                "grades": [18, 17, 20],
                "is_active": True
            }
            
            # Access and modify
            print(student["name"])
            student["age"] = 26
            student["major"] = "Computer Science"
            
            for key, value in student.items():
            print(f"{key}: {value}")
        \end{minted}
    \end{latin}
    \codecaption{کار با دیکشنری}
    \label{code:dictionary}
    
    %========================================
    % Chapter 6: Algorithms
    %========================================
    \chapter{الگوریتم‌ها}
    
    \section{جستجوی دودویی}
    
    \begin{latin}
        \begin{algorithm}
            \begin{algorithmic}[1]
                \Require Sorted array $A$, target $x$
                \Ensure Index of $x$ or $-1$
                \State $left \gets 0$, $right \gets n-1$
                \While{$left \leq right$}
                \State $mid \gets (left + right) // 2$
                \If{$A[mid] = x$}
                \State \Return $mid$
                \ElsIf{$A[mid] < x$}
                \State $left \gets mid + 1$
                \Else
                \State $right \gets mid - 1$
                \EndIf
                \EndWhile
                \State \Return $-1$
            \end{algorithmic}
        \end{algorithm}
        \algcaption[alg:binary]{جستجوی دودویی}
    \end{latin}
    
    \begin{theorem}
        جستجوی دودویی (\algref{alg:binary})
        در بدترین حالت $O(\log n)$ مقایسه انجام می‌دهد.
        \label{thm:binary}
    \end{theorem}
    
    %========================================
    % Chapter 7: Visualization
    %========================================
    \chapter{مصورسازی}
    
    \section{نمودار تابع}
    
    \begin{figure}[h]
        \centering
        \begin{latin}
            \begin{tikzpicture}
                \begin{axis}[
                    width=12cm, height=6cm,
                    xlabel={$x$},
                    ylabel={$f(x)$},
                    title={Common Functions},
                    grid=major,
                    legend pos=north west,
                    domain=-2:2,
                    samples=100
                    ]
                    \addplot[matinblue, thick] {x^2};
                    \addlegendentry{$x^2$}
                    \addplot[matinred, thick] {x^3};
                    \addlegendentry{$x^3$}
                \end{axis}
            \end{tikzpicture}
        \end{latin}
        \caption{نمودار $x^{2}$ و $x^{3}$}
        \label{fig:functions}
    \end{figure}
    
    \section{فلوچارت}
    
    \begin{figure}[h]
        \centering
        \begin{latin}
            \begin{tikzpicture}[
                node distance=2cm,
                arrow/.style={thick,->,>=stealth}
                ]
                \node[block] (start) {Start};
                \node[block, below of=start] (input) {Read Input};
                \node[decision, below of=input] (check) {Valid?};
                \node[block, right of=check, xshift=3cm] (process) {Process};
                \node[block, below of=check, yshift=-1.5cm] (error) {Error};
                \node[block, below of=process] (end) {End};
                
                \draw[arrow] (start) -- (input);
                \draw[arrow] (input) -- (check);
                \draw[arrow] (check) -- node[anchor=south] {Yes} (process);
                \draw[arrow] (check) -- node[anchor=east] {No} (error);
                \draw[arrow] (process) -- (end);
            \end{tikzpicture}
        \end{latin}
        \caption{فلوچارت ساده}
        \label{fig:flowchart}
    \end{figure}
    
    %========================================
    % CONCLUSION
    %========================================
    \chapter{جمع‌بندی}
    \label{chap:conclusion}
    
    این کتاب نقطه شروعی برای یادگیری برنامه‌نویسی
    و الگوریتم بود. مفاهیمی که آموختید:
    
    \begin{itemize}
        \item مبانی پایتون (\chref{chap:python} در کتاب اصلی)
        \item الگوریتم‌های جستجو و مرتب‌سازی
        \item ساختمان داده‌های پایه
        \item توابع و ماژول‌ها
    \end{itemize}
    
    برای ادامه مسیر، منابع زیر توصیه می‌شوند:
    \cite{knuth1984,cormen2009}.
    
    \vspace{1cm}
    \begin{flushleft}
        \textbf{پایان}\\
        تابستان ۱۴۰۵
    \end{flushleft}
    
    %========================================
    % BIBLIOGRAPHY
    %========================================
    \begin{latin}
        \printbibliography[title={منابع و مراجع}]
    \end{latin}
    
    %========================================
    % INDEX
    %========================================
    \printindex
    
\end{document}